From \parencite{wuest2016ml_manufacturing} :

"ML is already widely applied in different areas of manufacturing, e.g. optimization, control, and troubleshooting."
"These data-driven approaches are able to find highly complex and non-linear patterns in data of different types and sources and transform raw data to features spaces, so-called models, which are then applied for prediction, detection, classification, regression, or forecasting."
"Since most engineering and manufacturing problems are data-rich but knowledge-sparse, ML provides a tool to increase the understanding of the domain."
"ML has been successfully utilized in various process optimization, monitoring and control applications in manufacturing, and predictive maintenance in different industries."
"ML techniques were found to provide promising potential for improved quality control optimization in manufacturing systems, especially in 'complex manufacturing environments where detection of the causes of problems is difficult'."
"A very common challenge of ML application in manufacturing is the acquisition of relevant data. [...] the availability, quality, and composition [...] of the manufacturing data at hand have a strong influence on the performance of ML algorithms."
"In manufacturing practice, it is a common problem that values of certain attributes are not available or missing in the data-set."
"high-dimensional data can [...] contain a high degree of irrelevant and redundant information which may impact the performance of learning algorithms."
"In manufacturing application, supervised ML techniques are mostly applied due to the data-rich but knowledge-sparse nature of the problems [...] partly due to the availability of (a) expert feedback (e.g. quality) and (b) the labelled instances."
"often ML applications are found to be limited focusing on specific processes instead of the whole manufacturing program or manufacturing system."

From \parencite{weichert2019review_ml_optimization_production}

"The first step to solving a machine learning problem is to identify the data that is available. There are different types of data which will affect the type of analyses that can be used on them."
"The knowledge of the structure of production data is crucial for the right choice of machine learning models."
"Most data coming from production systems equipped with sensors can be considered as structured data which is much easier to handle than unstructured data."
"In a few cases, also nominal qualitative data are available [...] established and robust processes almost exclusively produce good parts and very few data are available for rejected parts. However, machine learning models need balanced data."
"Process optimization by machine learning can be structured into two main topics [...] optimization without change of production parameters during the manufacturing process [and] optimization with change of production parameters."
"Typical industrial applications for quality improvement based on machine learning are found in large-scale production such as plastic injection molding and the production of semiconductors [...] grounded on the high amounts of usable data points provided due to short cycle times."
"process complexity, the number of data points per input dimension to the model, and the model complexity are highly related."
"this connection was not respected, leading to complex models being trained on low amounts of data, risking overfitting and/or a lack of interpretability."
"The reviewed literature shows that the production-related applications of machine learning with or without optimization are manifold [...] there is hardly any correlation between the used data, the amount of data, the machine learning algorithms, the used optimizers, and the respective problem from the production."

From \parencite{plathottam2023review_ai_manufacturing} :

AI has the potential to be extremely helpful in manufacturing, particularly in applications such as predictive maintenance, quality assurance, and process optimization."
"The predictive and analytic capabilities of AI/ML models can provide useful insights to human operators to facilitate planning, support real-time decision-making, and improve manufacturing efficiency and safety."
"Predictive maintenance involves analyzing sensor data from equipment to anticipate potential equipment failures and scheduling maintenance routines to prevent unnecessary downtime [...] Typical manufacturing facilities average 15 h of downtime a week, with losses of about $20 000 per minute of production line downtime for large automotive companies."
"The increasing diversity and complexity of manufacturing tasks, workflows, and supply chains has increased the number of variables and interdependencies. Manually finding an optimal solution via experimentation is time- and resource-intensive."
"Getting the data needed to train AI models can be difficult for rare events or costly for large datasets that need labeling."
"Exposure to varying or extreme manufacturing environment conditions like heat release, robot motion, and pressure can affect sensor output, and over time, can lead to sensor drift which may bias collected data."
"AI/ML is often deployed in isolated, specialized situations, and thus does not pick up contextual information that could benefit the process further."
"The choice of which AI/ML solutions to implement for specific manufacturing problems is not always trivial and involves a certain degree of trade-offs."
"AI/ML implementation decisions must be suited to each company's unique situation and needs."

From \parencite{adekoya2022applications_ai_em} :

"Areas such as manufacturing, industrial design, inspection, monitoring and control, repairs and maintenance of industrial assets, product testing and evaluation have received their fair share [of AI application]."
"A very high percentage (77%) of ML application in EM was based on supervised learning, making SL technique common among machine learning techniques applied in EM."
"The high numbers of SL application in EM compared with other learning techniques can be partially attributed to the already labelled dataset's availability and SL lesser computational time than UL and RL."
"A high percentage of AI studies (28.9%) were applied in fault diagnosis than other domain areas in EM [...] followed by manufacturing monitoring, cost, and power consumption (14.5%)."
"There are no fast rules of superiority among ML algorithms. [...] their performance may differ from one area of application to the other."
"Several diverse theories, algorithms, and methods are available, which presents a challenge and a barrier in choosing the right AI technique for the appropriate engineering process or manufacturing process and environments."

From \parencite{presciuttini2024ml_iot_manufacturing_operations} :

"ML and DL algorithms can be the key to pursuing Operational Excellence and discovering hidden patterns in data collected through IoT sensors. These data can be used to develop different AI-based models to solve a wide range of tasks, from quality control-related problems to those concerning predictive maintenance."
"The Maintenance area is the most covered domain [...] most of the works cover Predictive Maintenance applications. Indeed, given the escalating demand to reduce downtime and economic losses, Predictive Maintenance is considered a prominent strategy for predicting anomalies in manufacturing systems."
"Quality is the second-most covered manufacturing operations area [...] IoT is considered pivotal for Quality Control, utilizing sensors for real-time defect detection and prompt issue rectification to prevent defective products from entering the market."
"Production is the most under-explored manufacturing operations area [...] the authors highlight that 75% of the potential research domains within the ML and Production Planning and Control intersection are either minimally explored or not addressed entirely."
"In the context of manufacturing applications, SL techniques are predominantly utilized because of the data-rich but knowledge-sparse nature of the tasks."
"The lack of interpretability of most ML and DL applications is one of the key barriers to their adoption in real industrial environments, leading to reduced trust in ML and DL predictions."

"ML can endow the PPC (Production Planning and Control) with the capacity of learning from historical or real-time data to react to predictable and unpredictable events."
"Artificial data is the most used data source in recent scientific literature. This probably highlights the difficulty of accessing data coming from companies."
"The first challenge is common when training ML models to identify disruptions. In fact, disruptive events in PPC (Production Planning and Control) such as machine breakdowns or quality problems tend to be scarce when compared to the total size of the dataset. Thus, ML techniques struggle to learn these events."
"Smart Planning and Scheduling is the most addressed use case in recent scientific literature [...] it normally uses structured data relatively easy to get from information systems."
"75% of the possible research domains in ML-PPC (Production Planning and Control) are barely explored or not addressed at all."
"Two key pitfalls are identified in the implementation of ML-PPC (Production Planning and Control) models: the complexity of using Internet of Things technologies to collect data and the difficulty of updating the ML model to adapt it to the manufacturing system changes."

From \parencite{panzer2022deep_reinforcement_learning_production_systems}:

"Deep RL differs in particular by the learning in direct interaction with its environment. It learns by a trial-and-error principle without requiring any pre-collected data or prior (human) knowledge and has the ability to adapt flexibly to uncertain conditions."
"Deep RL is applied in a wide range of production engineering disciplines [...] contributing to data-driven and flexible processes. In most applications, conventional methods were outperformed and implementation efforts or dependence on human experience were reduced."
"An application of deep RL was found in a wide range of production engineering disciplines" covering process control, scheduling, dispatching, assembly, maintenance, energy management, and quality control.
"In more than 85% of the total comparisons, deep RL algorithms outperformed the corresponding benchmarks and increased problem-specific performances."
"76% of the papers have validated the proposed solution within simulations. Only 24% of the papers conducted real-world testings [...] particularly high safety and reliability-related entry hurdles must be met, besides high system implementation efforts."

From \parencite{usuga2020ml_production_planning_control} :

"ML can endow the PPC with the capacity of learning from historical or real-time data to react to predictable and unpredictable events."
"Artificial data is the most used data source in recent scientific literature. This probably highlights the difficulty of accessing data coming from companies."
"The first challenge is common when training ML models to identify disruptions. In fact, disruptive events in PPC such as machine breakdowns or quality problems tend to be scarce when compared to the total size of the dataset. Thus, ML techniques struggle to learn these events."
"Smart Planning and Scheduling is the most addressed use case in recent scientific literature [...] it normally uses structured data relatively easy to get from information systems."
"75% of the possible research domains in ML-PPC are barely explored or not addressed at all."
"Two key pitfalls are identified in the implementation of ML-PPC models: the complexity of using Internet of Things technologies to collect data and the difficulty of updating the ML model to adapt it to the manufacturing system changes."

From \parencite{bergsto2026mlselection}:

"The largest share of articles is assigned to Phase 3 (Optimization), followed by Phase 2 (Development and Production). Phase 4 (Use and Reuse) represents a smaller portion of the dataset, while Phase 1 (Planning) contains the fewest articles."
"Supervised learning clearly dominates the literature across all product lifecycle phases and application areas. Classical supervised methods, including neural networks, random forests, and support vector machines, are most frequently reported."
"Across all clusters, supervised learning is the dominant paradigm, accounting for roughly 60–88% of the articles in each cluster. Reinforcement learning generally has the lowest share and counts, with a clear exception in Cluster 0 (control, process, scheduling), where it accounts for approximately 27%."
"Unsupervised learning shows more variation across clusters. It has the highest relative share in Cluster 4 (detection and anomaly-related applications), accounting for approximately 35% of the articles."
"The dominance of supervised learning suggests that method selection in industrial contexts is often shaped by data availability rather than by the problem's underlying structure."
"The widespread use of supervised learning should not automatically be seen as a weakness in industrial machine learning practice. In many production settings, reliability, clear validation, and predictable results are more important than the pursuit of new or experimental methods."
"The literature is strongly skewed toward the optimization phase of the product lifecycle, while early phases related to planning and concept development are comparatively underrepresented [...] optimization activities occur in data-rich operational settings, where structured process and sensor data, along with clearly defined performance objectives, support the application of machine learning methods."
"Many studies focus on isolated use cases and describe machine learning applications without clearly linking the problem being addressed, the available data, and the chosen methods."
"The most common methods are predominantly supervised learning approaches. Neural networks are the most frequently mentioned method, followed by random forest and support vector machines. Decision trees, convolutional neural networks, gradient boosting, and LSTM models also appear frequently."
"Across all phases, supervised learning is the dominant paradigm, accounting for approximately 73–80% of the articles in each phase. Unsupervised learning has its highest relative share in Phase 3 (Optimization), where it represents about 18% of the articles. Reinforcement learning remains the least represented paradigm across all phases, with relatively stable but low shares ranging from approximately 5% to 16%."
"With respect to RQ1, the results show a clear dominance of supervised learning across all lifecycle phases and application areas. This reflects the availability of labelled industrial data and the practical need for methods with clear evaluation metrics, while unsupervised and reinforcement learning remain less frequently reported."
"Regarding RQ2, machine learning research is strongly concentrated in later lifecycle phases, particularly optimization, whereas early phases such as planning and concept development are underrepresented. This suggests that current research activity is largely driven by data availability rather than by the potential long-term impact of early design decisions."
